\documentclass[10pt,a4paper]{amsart}
\usepackage[margin=19mm]{geometry}
\usepackage{amsmath,amssymb,mathtools}
\usepackage[scr=rsfs,cal=euler]{mathalfa}
\usepackage{xfrac}
\usepackage[unicode,hidelinks,bookmarks=false]{hyperref}
\newcommand{\ie}{i.e.\ }
\newcommand{\cf}{cf.\ }
\newcommand{\resp}{respectively\ }
\newcommand{\Subsection}{\subsection}
\providecommand{\nobreakdash}{\nobreak\hbox{-}}
\setlength{\emergencystretch}{2em}
\multlinegap0pt
\usepackage{xcolor}
\usepackage{tikz}
\newtheorem{theorem}{Theorem}
\newtheorem{proposition}{Proposition}
\newtheorem{lemma}{Lemma}
\title{Polyhedral norms and smooth Hahn-Banach extension}
\author{Saikat Roy}
\date{Source: arXiv:2606.03701v1 (CC BY 4.0)}
\begin{document}
\maketitle

\noindent\emph{Source and licence.} Polyhedral norms and smooth Hahn-Banach extension. Version arXiv:2606.03701v1 (CC BY 4.0), distributed by arXiv under the Creative Commons Attribution 4.0 International licence, http://creativecommons.org/licenses/by/4.0/, as recorded in the arXiv licence metadata for this exact version and shown on the abstract page https://arxiv.org/abs/2606.03701v1. Full licence notice: \texttt{licenses/arxiv-2606.03701-CC-BY-4.0.txt}.\par\smallskip

\noindent\textit{Editorial scope.} Counted unit: the article's theorem liftings (source0.tex lines 103--105). The article's complete original proof of it is delivered with it (source0.tex lines 175--293, one \texttt{proof} environment of 119 lines, which the article places away from the statement). No mathematics is written, repaired or re-derived: the statement and the proof are the article's own lines, in the article's own notation, and no retained line is substituted or re-flowed.  Retained uncounted as the source-local context the proof needs: lines 130--136; lines 149--151.  Nothing was omitted from any retained span, so the package carries no omission record.  No retained line contains a citation, so the wrapper carries no bibliography and no bibliography entry.  No retained line contains a figure environment, an image-inclusion command or drawing code, so no image is delivered and the package declares no asset dependency.  Editorial formatting only, in addition: an AMS \texttt{amsart} preamble that restates the article's own declarations and macros used by the retained lines, namely the environments \texttt{theorem}, \texttt{proposition}, \texttt{lemma} on separate counters, together with the standard packages those lines need.  The retained environments print in this document as Theorem 1, Proposition 1, Lemma 1 in order of appearance, whereas the article prints the same items with the numbering of its own full document.  The complete native source of the article as distributed in the arXiv e-print for this exact version is bundled unchanged as \texttt{sources/201994/source0.tex}.
\begin{proposition}\label{prop:smooth}
Let $X$ be a polyhedral space and $F$ be a face of $B_X$. Then,
\[
relint~(F)=\{\sum_{k=1}^n \lambda_k x_k:~~\sum_{k=1}^n \lambda_k=1, \lambda_k>0~\text{for all}~k\}.
\]
where $ext(F)=\{x_1,x_2,\dots,x_n\}$.
\end{proposition}

\begin{lemma}\label{smooth}
Let $Z$ be a finite-dimensional Banach space. Then, $f\in sm(B_{Z^*})$ if and only if $A_f=\{z_0\}$ for some unit vector $z_0$. 
\end{lemma}

\begin{theorem}\label{smooth HB liftings} 
Let $X$ be a polyhedral space and $Y$ be a subspace of it. Then any smooth functional $f$ on $Y$ extends to a norm-preserving smooth functional $\widetilde{f}$ on $X$ if and only if $f$ attains its absolute norm at an extreme point of $B_Y$ which is also an extreme point of $B_X$.
\end{theorem}

\begin{proof}[Proof of Theorem \ref{smooth HB liftings}]
First suppose that $f$ is a smooth functional on $Y$ that extends to a smooth functional  $\widetilde{f}$ on $X$ with the preservation of norm. Since $f$ is smooth, it attains its absolute norm at a unique point, i.e., $A_f=\{y_0\}$ for some unit vector $y_0$. Since  $\widetilde{f}$ is a Hahn-Banach extension of $f$, we have $A_f\subset A_{\widetilde{f}}$. At the same time $\widetilde{f}$ is a smooth linear functional, and thus it follows that $A_{\widetilde{f}}=\{y_0\}$. Therefore,  $\widetilde{f}$ is an exposing functional of $y_0$, and $y_0\in ext(B_X)$.


\medskip


\noindent Conversely, let $g$ be a smooth functional on $Y$ and $A_g=\{y_0\}$. Then $y_0\in ext(B_Y)$ and by the hypothesis $y_0\in ext(B_X)$. We now consider the (non-empty) collection
\[
S_{y_0}=\{\widetilde{f}\in ext(B_X):~\widetilde{f}(y_0)=1\}.
\]
Let $H$ be a subspace of $X$ defined as $H=\{x\in X:~\widetilde{f}(x)=0,~\widetilde{f}\in S_{y_0}\}$. Let $B_{y_0}=\{\widetilde{f}\in ext(B_{X^*}):~|\widetilde{f}(y_0)|<1\}$ and $C_{y_0}=\{\widetilde{f}\in ext(B_{X^*}):~|\widetilde{f}(y_0)|=1\}$. We observe that $C_{y_0}=S_{y_0}\cup (-S_{y_0})$ and $ext(B_{X^*})=B_{y_0}\cup C_{y_0}$. Also, the subspace $H$ is nothing but the intersection of kernels of all the linear functionals contained in $C_{y_0}$. If $H$ is a non-trivial subspace, we choose $u\in H$ and a real number $\lambda$ with sufficiently small absolute value such that $\max\{|\widetilde{f}(y_0\pm \lambda u)|:~\widetilde{f}\in B_{y_0}\}<1-\delta$ for some positive real number $\delta$. Also, it follows from the choice of $u$ that $|\widetilde{f}(y_0\pm \lambda u)|=1$ for each $\widetilde{f}\in C_{y_0}$. Thus,
\[
\max\{|\widetilde{f}(y_0\pm \lambda u)|:~\widetilde{f}\in ext(B_{X^*})\}=1.
\]
and we get $\|y_0\pm \lambda u\|=1$. However, then $y_0=\frac{1}{2}((y_0+\lambda u)+(y_0-\lambda u))$ is a convex combination of two unit vectors, a contradiction to the fact $y_0\in ext(B_X)$. This proves that $H$ must be trivial and shows $S_{y_0}$ contains a (algebraic) basis of $X^*$.


\medskip


\noindent Since $g\in sm(Y)$, $g\in relint~(W)$ for some facet $W$ of $B_{Y^*}$. Then by Proposition \ref{prop:smooth}
\[
g = \sum_{j=1}^p\lambda_jf_j,
\]
where 
\[
ext~(W)=\{f_1,f_2,\dots,f_p\}, \quad \sum_{j=1}^p\lambda_j=1, \quad \text{and}~ \lambda_j>0, ~\text{for all}~j.
\]
We now obtain a specific description of $W$. Let $\Pi:Y\to Y^{**}$ denote the canonical embedding. Then $\Pi(y_0)$ is an extreme point of $ext(B_{Y^{**}})$ and $\Pi(y_0)$ uniquely supports a facet of $B_{Y^*}$. Since $\Pi(y_0)(g)=1$, and $g\in relint~(W)$, we have $A_{\Pi(y_0)}=W$. Consequently, $f\in W$ if and only if $\Pi(y_0)(f)=f(y_0)=1$, which shows
\[
W=\{f\in B_{Y^*}:~f(y_0)=1\}.
\]
We now divide $S_{y_0}$ into two disjoint subsets
\begin{align*}
& S_{y_0}'=\{\widetilde{f}\in S_{y_0}:~\widetilde{f}|_Y=f,~f\in ext~(W)\},\\
& S_{y_0}''=\{\widetilde{f}\in S_{y_0}:~\widetilde{f}|_Y=f,~f\notin ext~(W)\}.
\end{align*}
Basically, $S_{y_0}'$ consists of members of $S_{X^*}$ that restricts to an extreme point of the facet $W$ of $B_{Y^*}$ and $S_{y_0}''$ represents the functionals on $X$ that restricts to an ordinary point (non-extreme) of the facet $W$.


\medskip



\noindent Let $I_j=\{\widetilde{f}\in S_{X^*}:~\widetilde{f}|_Y=f_j\}$ for each $j=1,2,\dots,p$. For convenience, assume
\[
I_j=\{\widetilde{f}_{1,j},\widetilde{f}_{2,j},\dots, \widetilde{f}_{n_j,j}\}, \qquad j=1,2,\dots,p,
\]
for natural numbers $n_1,n_2,\dots,n_p$. Let
\[
\mathcal{B}=(\bigcup_{j=1}^p I_j) \cup \mathcal{C},
\]
where $\mathcal{C}$ is a subset of $S_{y_0}''$ so that $\mathcal{B}$ contains a basis of $X^*$. Existence of such a set $\mathcal{B}$ is always guaranteed, since $S_{y_0}$ contains a basis of $X^*$. Anyway we may always choose $\mathcal{C} = \emptyset$, if the collection $(\bigcup_{j=1}^p I_j)$ already contains a basis of $X^*$. Let
\[
\mathcal{C}=\{\widetilde{g}_1,\widetilde{g}_2,\dots, \widetilde{g}_k\}~\qquad \text{and}  \qquad \widetilde{g}_i|_Y=g_i \quad \text{for each}~i=1,2,\dots,k.
\]
It now follows from the description of $W$ that $g_i\in W$ for each $i=1,2,\dots,k$.


\medskip



\noindent We now construct a Hahn-Banach extension $\widetilde{g}$ of $g$ such that $\widetilde{g}$ is a smooth linear functional on $X$. In particular, we shall show that the absolute norm attainment set of $g$ is a singleton set and use Lemma \ref{smooth} to conclude $\widetilde{g}$ is smooth. Let $\widetilde{g}:X\to \mathbb{R}$ be defined by
\[
\widetilde{g}(x)= \left[\sum_{j=1}^p\sum_{i=1}^{n_j}\alpha_j(\frac{1}{n_j}\widetilde{f}_{i,j})+\beta_{1}\widetilde{g}_{1}+\beta_{2}\widetilde{g}_{2}+\dots+\beta_{k}\widetilde{g}_k\right](x), \qquad x\in X,
\]


\noindent We now specify the choices for the scalars $\alpha_1, \alpha_2,\dots \alpha_p, \beta_{1}, \beta_{2}, \dots, \beta_k$:\\

\noindent Since $ext~(W)=\{f_1,f_2,\dots,f_p\}$, for each $r\in \{1,2,\dots, k\}$ there exist scalars $\mu_{1,r},\mu_{2,r}, \dots, \mu_{p,r}\in [0,1]$ with $\mu_{1,r}+\mu_{2,r}+ \dots +\mu_{p,r}=1$ such that
\[
g_{r}=\mu_{1,r}f_1+\mu_{2,r}f_2+ \dots +\mu_{p,r}f_p.
\]
For each $r$, choose $\beta_{r}>0$ such that
\[
\sum_{r=1}^{k}\beta_{r}\mu_{j,r}<\min\{\lambda_1,\lambda_2,\dots, \lambda_p\}, \quad \text{for each}~j\in \{1,2,\dots, p\}.
\]
Such a choice is always possible, since $\min\{\lambda_1,\lambda_2,\dots, \lambda_p\}>0$. Let
\[
\sum_{r=1}^{k}\beta_{r}\mu_{j,r}=t_j\lambda_j,\quad \text{for some}~t_j\in (0,1),~\text{for each}~j\in \{1,2,\dots, p\}.
\]
Also, let
\[
\alpha_j=(1-t_j)\lambda_j, \quad \text{for each}~j\in \{1,2,\dots, p\}.
\]
With such a specification, we claim that $\widetilde{g}$ is a smooth Hahn-Banach extension of $g$. Firstly, $\widetilde{g}$ is an extension of $g$. Observe that
\begin{align*}
\widetilde{g}\lvert_Y & = \left[\sum_{j=1}^p\sum_{i=1}^{n_j}\alpha_j(\frac{1}{n_j}\widetilde{f}_{i,j})+\beta_{1}\widetilde{g}_{1}+\beta_{2}\widetilde{g}_{2}+\dots+\beta_{k}\widetilde{g}_k\right]\Bigg\lvert_Y\\
& = \sum_{j=1}^p\alpha_jf_j+\beta_{1}g_{1}+\beta_{2}g_{2}+\dots+\beta_{k}g_k.
\end{align*}
It follows from the choice of the scalars $\alpha_1, \alpha_2,\dots \alpha_p, \beta_{1}, \beta_{2}, \dots, \beta_k$ that
\[
(\alpha_j+\sum_{r=1}^{k}\beta_{r}\mu_{j,r}) = (t_j\lambda_j+(1-t_j)\lambda_j) \qquad j\in \{1,2,\dots,p\}.
\]
Therefore, we get
\begin{align*}
\widetilde{g}\lvert_Y & = (\alpha_1+\sum_{r=1}^{k}\beta_{r}\mu_{1,r})f_1+(\alpha_2+\sum_{r=1}^{k}\beta_{r}\mu_{2,r})f_2+\dots+(\alpha_p+\sum_{r=1}^{k}\beta_{r}\mu_{p,r})f_p\\
& = (t_1\lambda_1+(1-t_1)\lambda_1)f_1+(t_2\lambda_2+(1-t_2)\lambda_2)f_2+\dots+(t_p\lambda_p+(1-t_p)\lambda_p)f_p\\
& = \lambda_1f_1+\lambda_2f_2+\dots+\lambda_pf_p\\
& = g.
\end{align*}

\medskip


\noindent Secondly, $\widetilde{g}$ is norm-preserving, since
\begin{align*}
1 =\|g\| & \leq \|\widehat{g}\|\\& \leq (\alpha_1+\sum_{r=1}^{k}\beta_{r}\mu_{1,r})+(\alpha_2+\sum_{r=1}^{k}\beta_{r}\mu_{2,r})+\dots+(\alpha_p+\sum_{r=1}^{k}\beta_{r}\mu_{p,r})\\
& =1.
\end{align*}
Thirdly, $\widetilde{g}$ attains its absolute norm only at $y_0$, since the collection $\mathcal{B}$ forms a basis of $X^*$, and
\[
A_{\widetilde{g}}=(\bigcap_{j=1}^{p}\bigcap_{i=1}^{n_j}A_{\widetilde{f}})\cap (\bigcap_{r=1}^{k}A_{\widetilde{g}_{r}}) = \bigcap_{l\in \mathcal{B}} (y_0+\ker~l)=\{y_0\}.
\]
Consequently, $\widetilde{g}$ is a smooth linear functional on $X$, and the proof is now complete.
\end{proof}

\end{document}
